\section{Circuits in place of programs}\label{sec:circuits}
\providecommand{\cG}{\mathcal G}\providecommand{\ctt}{\mathrm{tt}}\providecommand{\csize}{\mathrm{size}}
\providecommand{\ccore}{\mathrm{core}}\providecommand{\cdec}{\mathrm{Dec}}\providecommand{\cmux}{\mathrm{mux}}

H\"anni also asked for the same question with circuits in place of programs, and for an overall view of the
prevalence of malignity there. This section fixes one reading of that question, the \emph{circuit toy}, and
goes through the results of the note in it. The \emph{source} is the file
\texttt{research/circuits-final.md}
of the repository; it contains the proofs that are not given here. Results are theorems of the circuit toy
unless they are marked as computed, argued or guessed. The overall view is in \cref{sec:c-view}.

\subsection{The circuit toy}

\begin{definition}[Circuits]\label{def:c-circuit}
Fix $k\ge2$ input wires. A \emph{circuit of size $s$} is a list of gates $g_i=(\tau_i,a_i,b_i)$, $1\le i\le
s$,
with a truth table $\tau_i\in\B^4$ (all sixteen functions of two bits) and two ordered pointers
$a_i,b_i\in\{1,\dots,k+i-1\}$. Wires $1,\dots,k$ are the inputs and wire $k+i$ is gate $i$, whose value is
$\tau_i$ applied to the values of the wires $a_i$ and $b_i$. The output is the last gate. $\cG_{k,s}$ is the
set of these circuits, $|\Gamma|$ is the size of a circuit $\Gamma$, $f_\Gamma\colon\B^k\to\B$ is the
function
it computes, and
\[
N_s=|\cG_{k,s}|=\prod_{i=1}^{s}16\,(k+i-1)^2,\qquad p_i=\frac1{16\,(k+i-1)^2}.
\]
A gate is \emph{active} if it is the last gate or an ancestor of it, and \emph{dead} otherwise. A circuit is
\emph{reduced} if all its gates are active. The \emph{core} of $\Gamma$ is the reduced circuit obtained by
deleting the dead gates and renumbering in order; it computes $f_\Gamma$. For a function $f$, $\csize(f)$ is
the least size of a circuit that computes it.
\end{definition}

\paragraph{Priors.} The \emph{fixed-size prior} $\mu_s$ is uniform on $\cG_{k,s}$. A \emph{mixture prior}
has
weights $w_s>0$ for $1\le s\le s_{\max}$, with $s_{\max}\le\infty$ and $\sum_sw_s=1$, and gives
$\Gamma\in\cG_{k,s}$ the weight $q(\Gamma)=w_s/N_s$; we put $\varrho(\sigma,a)=w_{\sigma+a}/w_\sigma$.
Assumption (A1), used only where named, is that $w_s$ is nonincreasing. The \emph{length prior} is the
mixture
with $w_s=2^{-s}$ and $s_{\max}=\infty$: the size in unary, then the gates. It satisfies (A1), and
$\varrho(\sigma,a)=2^{-a}$. These are the \emph{priors of the family}, and $\mu$ stands for any of them. The
\emph{core weight} is
\[
\beta_s(a)=p_s\,e_{a-1}(p_1,\dots,p_{s-1})\ \ (a\le s),\qquad\beta_s(a)=0\ \ (a>s),\qquad
\beta_q(a)=\sum_sw_s\beta_s(a),
\]
with $e_i$ the elementary symmetric polynomial of degree $i$; $\beta$ stands for $\beta_s$ or $\beta_q$.

\paragraph{Hypotheses, data, queries.} A hypothesis is a circuit, read as the function it computes.
\emph{Data} are $D=(I,f^*|_I)$ for a set $I\subseteq\B^k$ of inputs and a target $f^*$. A circuit $\Gamma$
\emph{fits} $D$ if $f_\Gamma=f^*$ on $I$. Let $\xi(D)$ be the prior probability that a circuit fits $D$,
$\ell_D=-\log_2\xi(D)$, and $s_D$ the least size of a fitting circuit. For a query $u\notin I$, or a tuple
of
queries, the predictor's distribution is $\pi(y\mid u,D)=\mu\{f_\Gamma(u)=y\mid\Gamma\text{ fits }D\}$, and
the
sampling predictor returns $Y\sim\pi(\cdot\mid u,D)$.

In the \emph{sequence form}, or \emph{index reading}, the $k$ inputs carry the binary digits
$b_0,\dots,b_{k-1}$ of $t-1$ for a position $t\in\{1,\dots,2^k\}$, and the circuit says that the bit at
position $t$ is $f_\Gamma(t)$. Its \emph{truth table} is $\ctt(\Gamma)=f_\Gamma(1)\cdots f_\Gamma(2^k)$, and
$\Gamma$ \emph{prints} $x$ if $x\pre\ctt(\Gamma)$. A context $x$ of length $n$ is the data set with
$I=\{1,\dots,n\}$ and labels $x$, and a block of $m$ bits answers the queries $n+1,\dots,n+m$. Then
\[
\xi(x)=\mu\{\Gamma:\ x\pre\ctt(\Gamma)\},\qquad\pi_m(y\mid x)=\frac{\xi(xy)}{\xi(x)}\quad(|x|+m\le2^k),
\]
and $\csize(x)$ is the least size of a circuit that prints $x$. The sequence form is the special case of the
query form in which the data are an initial segment, and every result for data and queries applies to it. At
a cut $n=2^j$ the index bits $b_j,\dots,b_{k-1}$ are $0$ on the whole context, and positions $n+1,\dots,2n$
differ from positions $1,\dots,n$ only in $b_j$. So a dyadic cut is a query that differs from the data in a
bit that never varied.

We use the index reading because two other natural readings, with one gate or one pointer per output bit,
give a predictor that learns nothing, or only a frequency.

\begin{theorem}[Two degenerate readings]\label{thm:c-readings}
\begin{enumerate}[label=(\alph*),leftmargin=2em,itemsep=2pt,topsep=2pt]
\item \emph{One gate per output bit.} Let a circuit with $r\ge1$ random inputs be uniform on $\cG_{r,s}$,
$s\ge N$, and let its $N$ output bits be the values of its last $N$ gates. Then for every fixed value of the
inputs the output is uniform on $\B^N$. So the predictor has $\pi_m(y\mid x)=2^{-m}$ for every context and
block, also for mixtures over $r$ and $s\ge N$.
\item \emph{One pointer per output bit.} Let the prior on $r$, $s$ and the gates be arbitrary and let, given
these, the $N$ output pointers be independent and uniform on the $r+s$ wires. Let $\phi$ be the fraction of
wires that carry $1$, a function of the gates and the inputs. Then the probability $\xi(x)$ that the output
begins with $x\in\B^n$, a string with $i$ ones, is $\mathbb E[\phi^{\,i}(1-\phi)^{\,n-i}]$. So $\xi$ is
exchangeable, and $\pi_m(y\mid x)$ depends on $x$ and $y$ only through their numbers of ones.
\end{enumerate}
\end{theorem}
\begin{proof}
(a) Condition on the inputs and on the gates $1,\dots,s-N$. For $i=s-N+1,\dots,s$ in turn, the value of gate
$i$ is the entry of $\tau_i$ at the pair of values of the wires $a_i,b_i$. That pair is determined by the
inputs, the earlier gates and $(a_i,b_i)$, all independent of $\tau_i$, whose four entries are independent
fair bits. So the value is a fair bit independent of everything earlier. (b) Given the gates and the inputs,
the output bits are independent draws of the value of a uniform wire.
\end{proof}

The theorem is about these two readings. It does not show that every reading with one wire per output bit is
degenerate, nor that the index reading is the only sensible one. In the index reading the link between
context
and continuation is the reuse of the same gates at different positions. A circuit that takes the context as
input and is applied autoregressively is a third reading; it reuses a description across steps, belongs with
\cref{prop:c-recurrence}, and is not treated here.

\paragraph{Worlds, critical uses, takeover.} A \emph{world} is a circuit $W$, with target $f_W$. It is
\emph{in the class} if $\mu\{\Gamma:f_\Gamma=f_W\}>0$. A \emph{critical use} is a triple of a world, data
and a
query, fixed from outside: the world hands the predictor $D$, receives $Y$ at the query, and acts on it. Its
\emph{own continuation} $d$ is $f_W$ at the query, and the \emph{deviation mass} is $1-\pi(d\mid u,D)$. In
the sequence form a \emph{call rule} $\mathcal C$ is a set of pairs $(W,n_W)$, at most one for each circuit;
with $x_W=\ctt(W)_{1:n_W}$ and a mixture prior $q$, the quantities of \cref{def:call} are
\[
C(x)=\sum_{W:\,x_W=x}q(W),\qquad H_{\mathcal C}=\sum_{(W,n_W)\in\mathcal C}q(W),\qquad
\kappa(x)=\frac{C(x)}{\xi(x)}.
\]
A \emph{test} is a set $G$ of nonempty strings, with first hits as in \cref{def:world}; under a test every
circuit that prints a first hit hands over there, and $\kappa=1$. The forms of \cref{def:mph} are taken
over:
the majority form $\pi(d\mid u,D)<\frac12$, the payload form $\pi(A\mid u,D)\ge\theta$ for a set $A\not\ni
d$,
and the critical-use form $\Pr[Y\in A(X)]\ge\theta$ over a population of uses. As in the note they are
necessary for calling the prior malign at a world and not sufficient.

A takeover needs a decoder. In the note the block $y$ is run as $U$ on $y0^\infty$. Here hypotheses are
circuits and a block is a string, so the world applies a fixed decoding map to it. The \emph{literal code}
of
$\Gamma\in\cG_{k,s}$ is $e(\Gamma)=\bar e(s)\,\tau_1[a_1][b_1]\cdots\tau_s[a_s][b_s]$, with pointer fields
of
$\lceil\log_2(k+s)\rceil$ bits; a field with value $c$ at gate $i$ is read as wire $1+(c\bmod(k+i-1))$. Its
length is $l(s)=|\bar e(s)|+s\,(4+2\lceil\log_2(k+s)\rceil)$. $\cdec(y)$ is the circuit decoded from the
front
of $y$ if the rest of $y$ is $10^i$ for some $i\ge0$. The counterpart of \cref{def:sr} is
\[
\SR_m(x)=\{y\in\B^m:\ \cdec(y)\text{ is defined and }xy\pre\ctt(\cdec(y))\}.
\]
The decoder is an interpreter that lives outside the hypothesis class.

\paragraph{What the class assumes.} These are decisions, not results. (i) Plain circuits: fan-in two, each
gate evaluated once per output, no gate used twice, no depth bound, no weight sharing. (ii) Hypotheses are
total: every circuit has a bit at every position, so the worlds of the note that hand over and then print
nothing, and the quantity $\varphi$ of \cref{prop:honest}, cannot be expressed; this favours the honest
side.
(iii) The class depends on the horizon through $k$, and which positions are cheap to name depends on the
binary coding of the index. (iv) The prior is on ordered gate lists, with dead gates allowed. (v) The full
binary basis, with uniformly drawn truth tables. (vi) The sampling predictor. (vii) The world is in the
predictor's class wherever a bound compares with the world. (viii) The use is fixed from outside, and a
hypothesis can recognise it only through its own input, the index or the query.

\subsection{What transfers unchanged}

\begin{theorem}[Measure, calibration, budget, feedback]\label{thm:c-transfer}
Let $\mu$ be a prior of the family and write $\mu(f)=\mu\{\Gamma:f_\Gamma=f\}$.
\begin{enumerate}[label=(\alph*),leftmargin=2em,itemsep=2pt,topsep=2pt]
\item \emph{Measure.} In the sequence form $\xi(\eps)=1$ and $\xi(x0)+\xi(x1)=\xi(x)$ for $|x|<2^k$. Hence
$Z_m(x)=\sum_{y\in\B^m}\xi(xy)=\xi(x)$ for $|x|+m\le2^k$.
\item \emph{Calibration.} Fix $I$ and a query $u\notin I$. Let $\Gamma^*\sim\mu$, $D=(I,f_{\Gamma^*}|_I)$,
and
let $Y\sim\pi(\cdot\mid u,D)$ be drawn independently of $\Gamma^*$ given $D$. Then given $D$, the answer
$f_{\Gamma^*}(u)$ and $Y$ are independent with the common law $\pi(\cdot\mid u,D)$. In the sequence form let
$G$
be a test, let the world be drawn from $\mu$, and let $X$ be the first hit of $G$ among the prefixes of its
truth table, on the event that there is one of length at most $2^k-m$. Given $X=x$, the world's next $m$
bits
and $Y\sim\pi_m(\cdot\mid x)$ are independent with the law $\pi_m(\cdot\mid x)$.
\item \emph{Budget.} Let $z=\ctt(W)$ with $\xi(z)>0$. Then
$\sum_{n=0}^{2^k-m}\bigl(1-\pi_m(z_{n+1:n+m}\mid z_{1:n})\bigr)\le m\ln(1/\xi(z))$, and the term is at least
$\frac12$ at no more than $m\log_2(1/\xi(z))$ cuts $n$. For a mixture prior, $\xi(z)\ge w_\sigma/N_\sigma$
if a
circuit of size $\sigma\le s_{\max}$ prints $z$. Let queries $u_1,u_2,\dots$ be answered one at a time,
repetitions allowed, each depending on everything revealed before, the true answer being added to the data
after each, which gives $D_t$. If $\mu(f^*)>0$ and $a_t=\pi(f^*(u_t)\mid u_t,D_{t-1})$, then
$\sum_t(1-a_t)\le\ln(1/\mu(f^*))$, and $a_t\le\frac12$ for at most $\log_2(1/\mu(f^*))$ values of $t$.
\item \emph{One query off the data.} For $u^\dagger\notin I$ and every $f'$ that fits $D$ with
$f'(u^\dagger)\ne f^*(u^\dagger)$:
$\ \mu(f^*)/\xi(D)\le\pi(f^*(u^\dagger)\mid u^\dagger,D)\le1-\mu(f')/\xi(D)$.
\item \emph{Feedback.} Sequence form, with a response map $R$ as in \cref{sec:feedback}: the predictor
returns
$Y\sim\pi_m(\cdot\mid x)$ and the record continues with $R(Y)$. The prediction comes true with probability
$\pi_m(\Fix(R)\mid x)$. If $Y=y\in\Fix(R)$, no circuit that prints $xd$ with $d\ne y$ prints any
$x'\sqsupseteq xy$. If $R$ is the identity in every round, the record continues with $u\in\B^{jm}$, where
$|x|+jm\le2^k$, with probability exactly $\xi(xu)/\xi(x)$; it continues with a block $d$ such that no
circuit
in the support prints $xd$ with probability $0$. If the predictor returns the more probable next bit, $0$ on
a tie, and the record continues with the other bit from position $n$ on, then
$\xi(x_{1:n+j})\le2^{-j}\xi(x_{1:n})$ while $n+j\le2^k$; for a mixture prior with $s_{\max}<\infty$ and
$q_{\min}=\min_sw_s/N_s$, $\xi(x_{1:n+j})=0$ for all $j>\log_2(1/q_{\min})$ with $n+j\le2^k$.
\end{enumerate}
\end{theorem}
\begin{proof}
(a) Every circuit has a truth table of exactly $2^k$ bits, so the sets $\{\Gamma:x\pre\ctt(\Gamma)\}$,
$x\in\B^n$, partition the circuits. (b) is proved as \cref{thm:calibration}:
$\Pr[f_{\Gamma^*}(u)=y\mid D]=\pi(y\mid u,D)$ by definition; first hits are pairwise incomparable, so $X=x$
with
next block $d$ has probability $\xi(xd)$, and $X=x$ has probability $\xi(x)$ by (a). (c) is the proof of
\cref{thm:budget}(iv), which uses only that the conditional probabilities of the true bits multiply to
$\xi(z)$; and $\xi(z)$ is at least the weight of one circuit that prints $z$. For queries,
$a_t=\xi(D_t)/\xi(D_{t-1})$ telescopes, and every circuit computing $f^*$ fits every $D_t$. (d) The circuits
computing $f^*$, respectively $f'$, are among those fitting $D$. (e) The proofs of
\cref{thm:feedback,thm:contrarian} use only (a) and that a positive value of $\xi$ is at least the weight of
one circuit; by (a) the product in \cref{thm:feedback}(c) telescopes to an equality.
\end{proof}

By (a) the first qualification of \cref{thm:calibration} disappears: no world prints fewer than $m$ further
bits. The hypotheses of (b) remain. The world is drawn from the predictor's own prior, and the use is read
off
the data, by a test or by a query chosen from the data alone. With a call rule that selects users by
something
the data do not show, calibration does not apply. The factor $m$ in the first budget and its absence in the
second belong to two protocols: overlapping blocks at every cut, and answers revealed whole. Parts (b) and
(c)
do not bound the deviation at one query that something singled out. A use is such a point, and (d) is what
is
proved about it. In (e) the ``dark room'' of \cref{sec:feedback} becomes: an obedient world follows a sample
from the circuit prior conditioned on the past, and is moved off any default that no circuit in the support
prints. For a world outside the class nothing is guaranteed. The hardest function of $k$ inputs needs about
$2^k/k$ gates~\cite{frandsen05}, and fewer than $2N_s$ functions have circuits with at most $s$ gates, so
for
arbitrary sequences of length $2^k$ this is the generic case unless the size bound is exponential in $k$.

\subsection{The core law and the price of an active gate}

\begin{theorem}[Core law; price of an active gate]\label{thm:c-core}
\begin{enumerate}[label=(\alph*),leftmargin=2em,itemsep=2pt,topsep=2pt]
\item For a reduced circuit $\chi$ of size $a$ and $s\ge a$: $\
\mu_s\{\Gamma:\ccore(\Gamma)=\chi\}=\beta_s(a)$,
and for a mixture prior $q\{\Gamma:\ccore(\Gamma)=\chi\}=\beta_q(a)$. So under every prior of the family the
mass
of a core depends only on its size, the posterior on cores is $\beta(|\chi|)\mathbf 1[\chi\text{ fits }D]$
normalised, and every prediction is a function of the core. Consequently, for $a=\csize(f)\le s$,
\[
\mu_s\{\Gamma:f_\Gamma=f\}\ \ge\ \beta_s(a)\ \ge\ \frac{p_s}{N_{a-1}},\qquad\text{and under the length
prior}\quad
q\{\Gamma:f_\Gamma=f\}\ \ge\ \frac{2^{-a}}{N_a}.
\]
\item For $1\le a\le s-1$,
\[
\frac1{16a}\Bigl(\frac1{k+a-1}-\frac1{k+s-1}\Bigr)\ \le\ \frac{\beta_s(a+1)}{\beta_s(a)}\ \le\
\frac1{16a}\sum_{j=k}^{k+s-2}\frac1{j^2}\ \le\ \frac1{16a(k-1)}.
\]
The lower bound is at least $\frac1{32a(k+a-1)}$ when $s\ge2a+k-1$. For every prior of the family
$\beta(a+1)\le\beta(a)/(16a(k-1))$. So $\beta$ is decreasing, and
$\beta(a)\le\frac1{16k^2}\,(16(k-1))^{-(a-1)}/(a-1)!$ for $a\ge1$.
\end{enumerate}
\end{theorem}
\begin{proof}
(a) A circuit $\Gamma\in\cG_{k,s}$ with core $\chi$ is determined by the set $J\subseteq\{1,\dots,s\}$ of
its
active positions, with $|J|=a$ and $s\in J$, and by its dead gates: the active gate at the $j$-th position
of
$J$ has the truth table of the $j$-th gate of $\chi$ and its pointers, renumbered. Conversely, for every
such
$J$ and every choice of the gates at the positions $i\notin J$, of which there are $1/p_i$ at position $i$,
the resulting list is a circuit whose active set is exactly $J$: the gates at $J$ read only inputs and gates
at
$J$, all of them are ancestors of gate $s$ because $\chi$ is reduced, and no gate outside $J$ is read by a
gate in $J$. So the number of circuits with core $\chi$ is
$\sum_J\prod_{i\notin J}p_i^{-1}=N_s\sum_J\prod_{i\in J}p_i=N_s\,p_s\,e_{a-1}(p_1,\dots,p_{s-1})$. For a
mixture,
sum over $s$ with the weights $w_s$. For the inequalities: a minimum circuit for $f$ is reduced;
$e_{a-1}(p_1,\dots,p_{s-1})\ge p_1\cdots p_{a-1}=1/N_{a-1}$; and the term $s=a$ of $\beta_q(a)$ is
$2^{-a}/N_a$ for the length prior.

(b) Write $e_i=e_i(p_1,\dots,p_{s-1})$, so that the ratio is $e_a/e_{a-1}$. Counting pairs of an $(a-1)$-set
$J$
and an index $i\notin J$ gives $a\,e_a=\sum_{|J|=a-1}\prod_{i\in J}p_i\cdot\sum_{i\notin J}p_i$. The $p_i$
decrease, so $\sum_{i\notin J}p_i\ge\sum_{i=a}^{s-1}p_i\ge\frac1{16}\int_{k+a-1}^{k+s-1}x^{-2}dx$, which is
the
lower bound; and $\sum_{i\notin J}p_i\le
e_1=\frac1{16}\sum_{j=k}^{k+s-2}j^{-2}\le\frac1{16}\int_{k-1}^\infty
x^{-2}dx$, which gives the upper bounds. $\frac1{k+a-1}-\frac1{k+s-1}\ge\frac1{2(k+a-1)}$ exactly when
$s\ge2a+k-1$. For a mixture apply the upper bound to each $s\ge a+1$ and sum with the weights $w_s$. Finally
$\beta_s(1)=p_s\le p_1=1/(16k^2)$, also after mixing, and the ratio bound is applied $a-1$ times.
\end{proof}

So every prior of the family is a prior on cores, and a function with a small circuit keeps mass at least
$\beta(\csize(f))$ however large the nominal size. One more active gate costs at least $\log_2(16a(k-1))$
bits
under every prior of the family, and at most $\log_2(32a(k+a-1))$ bits under $\mu_s$ with slack
$s\ge2a+k-1$. Under the length prior no upper bound of that kind is proved, and the computed price is
higher:
$13.00$ bits at $k=20$, $a=1$ and $14.12$ at $a=8$, against the lower bounds $8.25$ and $11.25$. The theorem
does not give the mass of a function, because the number of cores per function is not controlled: $1/N_s$ is
a lower bound on a mass and not an estimate. Under $\mu_s$ the active part of a typical circuit has of order
$\sqrt{k+s}$ gates. This follows from a limit law of Janson for random recursive directed acyclic graphs, as
stated in~\cite{burghart26}; the identification of a uniform gate list with that model is ours. So ``a
circuit of size $s$ does $s$ operations'' is true of the list and not of the computation.

\subsection{Triggers and the floor}

A hypothesis that explains the data and then gives a chosen answer at a use must recognise the use from its
own input. Let $z(n)$ be the number of trailing zeros of $n$ in binary.

\begin{lemma}[Sizes of triggers and payloads]\label{lem:c-trigger}
Over the full binary basis, on the index bits $b_0,\dots,b_{k-1}$ of $t-1$:
\begin{enumerate}[label=(\alph*),leftmargin=2em,itemsep=2pt,topsep=2pt]
\item \emph{Position test.} For $1\le n<2^k$ the function $[\,t-1\ge n\,]$, which is $0$ on positions
$1,\dots,n$ and $1$ afterwards, is computed with $k-1-z(n)$ gates; with $0$ gates it is the index wire
$b_{k-1}$. So at the $2^{g+1}-1$ multiples $n$ of $2^{k-1-g}$ below $2^k$ the test needs at most $g$ gates.
\item \emph{Dyadic cut.} If $n=2^j$ and $m\le2^j$, the index wire $b_j$ is $0$ on positions $1,\dots,n$ and
$1$
on positions $n+1,\dots,n+m$: the test costs no gate.
\item \emph{Payload.} Every Boolean function of $j\ge2$ variables has a circuit of at most $2^j-3$ gates.
Hence for every block $y\in\B^m$ with $n+m\le2^k$ there is a circuit $P_{n,y}$ on the index with
$P_{n,y}(n+i)=y_i$ for $1\le i\le m$ and at most $\max(1,2^{\lceil\log_2m\rceil}-3)\le\max(1,2m-5)$ gates.
\item \emph{Selector.} $\cmux(c;u,v)=(c\wedge v)\vee(\neg c\wedge u)$ is three gates. A constant payload
needs
one gate: $u\vee c$ for the value $1$ and $u\wedge\neg c$ for the value $0$.
\end{enumerate}
\end{lemma}
\begin{proof}
(a) Let $c_i$ be the digits of $n$ and $G_i=[\,(b_i\dots b_0)\ge(c_i\dots c_0)\,]$. $G_0$ is $b_0$ if
$c_0=1$
and the constant $1$ if $c_0=0$, and $G_i=b_i\wedge G_{i-1}$ if $c_i=1$, $G_i=b_i\vee G_{i-1}$ if $c_i=0$.
With
$j=z(n)$: $G_i$ is the constant $1$ for $i<j$, $G_j=b_j$ is an index wire, and each step $i=j+1,\dots,k-1$
adds one gate. (b) Positions $1,\dots,2^j$ have $t-1<2^j$, and positions $2^j+1,\dots,2^{j+1}$ have $b_j=1$.
(c) A function of two variables is one gate, and $f=\cmux(x_j;f_0,f_1)$ gives $M(j)\le2M(j-1)+3$ for the
largest size $M(j)$ needed, so $M(j)+3\le2^{j-2}(M(2)+3)=2^j$. For the block let $j=\lceil\log_2m\rceil$.
The
$m\le2^j$ numbers $t-1\in\{n,\dots,n+m-1\}$ are distinct modulo $2^j$, so a function of the low $j$ index
bits
has the required values; for $j\le1$ one gate suffices. (d) By inspection.
\end{proof}

The count in (a) is exact (source, Lemma~7(a)). A test on the data is dearer in the index reading, because
each position is computed from scratch: to make a one-bit payload depend on $w$ earlier bits of the stream
of
an explanation $E$ of $\sigma$ gates through a test of $g_\phi$ gates, the source uses $w$ shifted copies of
$E$ and at most $w(\sigma+2k)+g_\phi+k+1$ further gates (Lemma~7(f) there). That nothing essentially cheaper
exists is not proved.

\begin{theorem}[Floor for a triggered payload]\label{thm:c-floor}
Let $\vartheta$ be a circuit of $g\ge0$ gates on the inputs that is $0$ on $I$ (for $g=0$, an input wire),
and
$P$ a circuit of $|P|$ gates on the inputs. For $E\in\cG_{k,\sigma}$ let $\mathrm{Tw}[E]$ be $E$ followed by
the gates of $\vartheta$, the gates of $P$, and $\cmux(\vartheta;E,P)$. It has $\sigma+a$ gates,
$a=g+|P|+3$,
fits $D$ exactly when $E$ does, and equals $P$ wherever $\vartheta=1$. For a constant payload take $a=g+1$.
\begin{enumerate}[label=(\alph*),leftmargin=2em,itemsep=2pt,topsep=2pt]
\item \emph{Mixture prior.} $\ q\{\Gamma:\ \Gamma\text{ fits }D,\ f_\Gamma=P\text{ on }\{\vartheta=1\}\}\
\ge\
\sum_{\sigma\le s_{\max}-a}w_{\sigma+a}\,(N_\sigma/N_{\sigma+a})\,\mu_\sigma\{\Gamma\text{ fits }D\}$. If a
share
$\rho_{\sigma_0}$ of the posterior is on circuits of size at most $\sigma_0\le s_{\max}-a$, the posterior
share
of the payload on $\{\vartheta=1\}$ is at least
$\rho_{\sigma_0}\cdot\min_{\sigma\le\sigma_0}\varrho(\sigma,a)\cdot\bigl(16\,(k+\sigma_0+a-1)^2\bigr)^{-a}$.
\item \emph{Length prior.} For $\epsilon\in(0,1)$ and $s_0=\lceil\ell_D+\log_2(1/\epsilon)\rceil$ the share
is
at least $(1-\epsilon)\bigl(32\,(k+s_0+a-1)^2\bigr)^{-a}$, and $\ell_D\le s_D+\log_2N_{s_D}$.
\item \emph{Critical uses.} For a call rule and blocks $y_x$ with wrappers of $a_x$ gates, $\Pr[Y=y_X]$ is
at
least the expectation over uses of the bound in (a). If the calling circuits with context $x$ have size at
most $s_{\max}-a_x$, the wrappers around them alone have share at least
$\min_\sigma\bigl[\varrho(\sigma,a_x)N_\sigma/N_{\sigma+a_x}\bigr]\cdot\kappa(x)$ at $x$.
\item \emph{A never-varied bit.} Let the prior be the length prior, let an input coordinate $v$ be $0$ on
all
of $I$ and $1$ at the query $u^\dagger$, and $s_0=\lceil\ell_D\rceil+1$. Then
$\pi(y\mid u^\dagger,D)\ge1/(32(k+s_0)^2+2)$ for each $y\in\B$, and for each $y$ the fitting circuits that
answer $y$ at every input with $u_v=1$ have posterior share at least $1/(32(k+s_0)^2)$. In the sequence
form,
at a cut $n=2^j$ with $0\le j<k$, each value of the next bit has share at least $1/(32(k+s_0)^2+2)$ and each
of the two constant continuations of any length $m\le2^j$ has share at least $1/(32(k+s_0)^2)$, where
$s_0=\lceil-\log_2\xi(x)\rceil+1\le\csize(x)+\log_2N_{\csize(x)}+2$.
\end{enumerate}
\end{theorem}
\begin{proof}
(a) $E\mapsto\mathrm{Tw}[E]$ is injective from $\cG_{k,\sigma}$ into $\cG_{k,\sigma+a}$, because $E$ is the
list of the first $\sigma$ gates. It maps fitting circuits to fitting circuits with the payload on
$\{\vartheta=1\}$. So the fraction of $\cG_{k,\sigma+a}$ with these properties is at least
$(N_\sigma/N_{\sigma+a})\mu_\sigma\{\Gamma\text{ fits }D\}$; multiply by $w_{\sigma+a}$ and sum. For
$\sigma\le\sigma_0$, $N_{\sigma+a}/N_\sigma\le(16(k+\sigma_0+a-1)^2)^a$. (b) The length prior gives the
circuits
larger than $s_0$ the mass $2^{-s_0}$, so their posterior mass is at most $2^{\ell_D-s_0}\le\epsilon$; and
$\varrho=2^{-a}$. One fitting circuit of size $s_D$ has weight $2^{-s_D}/N_{s_D}$. (c) Average (a) over the
contexts; for the second claim restrict the injection to calling circuits. (d) The posterior mass of the
circuits larger than $s_0$ is at most $\frac12$. For $\Gamma\in\cG_{k,s}$ let $\Gamma^\oplus$ and
$\Gamma^{\oplus\prime}$ be $\Gamma$ with one more gate, the exclusive or of gate $s$ and $v$ in the two
pointer
orders. Both compute $f_\Gamma(u)\oplus u_v$, fit $D$ exactly when $\Gamma$ does, and flip the answer at
$u^\dagger$. They are two distinct circuits that determine $\Gamma$, and their weights sum to
$q(\Gamma)\,p_{s+1}$. Let $\eta=p_{s_0+1}=1/(16(k+s_0)^2)$ and $\bar y=1-y$. Summing over the fitting
$\Gamma$ of
size at most $s_0$ with answer $y$ gives $\pi(\bar y)\ge\eta\,(\pi(y)-\frac12)=\eta\,(\frac12-\pi(\bar y))$,
so
$\pi(\bar y)\ge\eta/(2(1+\eta))$. For the second claim append instead the gate of
\cref{lem:c-trigger}(d) for the constant $y$, in both pointer orders, and sum over all fitting $\Gamma$ of
size at most $s_0$: the share is at least $\eta/2$. The sequence form follows from \cref{lem:c-trigger}(b).
\end{proof}

Three differences from \cref{thm:twin}. First, the cost is not a constant. A gate appended to an explanation
of $\sigma$ gates costs at most $4+2\log_2(k+\sigma+a-1)+\log_2(1/\varrho)$ bits, which grows with the
logarithm of the size of the explanation, though not with the amount of data beyond that. Second, the test
is
paid in gates. The counterpart of the $K(n)$ of the note is $k-1-z(n)$ gates; like $K(n)$ it is small at a
few positions, and it is zero at every dyadic cut with a block no longer than the context. Third, a test on
the data costs copies of the explanation; a hypothesis that received the context as input would pay for the
test alone. In numbers, for $k=20$ and explanations of $1000$ gates an appended gate costs about $24.0$
bits.
A one-bit deviation costs about $24$ bits at a dyadic cut and at most $480$ bits at the positions hardest to
name; a block of $64$ arbitrary bits adds at most $64$ gates for payload and selector, about $1540$ bits. In
the program toy the same wrappers cost $K(n)+O(1)$ and $K(n)+64+O(1)$ bits. In (d), for $k=20$ and
$s_D=1000$, $s_0\approx22\,326$ and the bound is $2^{-33.9}$.

Parts (a) to (c) are floors for one construction and not estimates of the deviation mass: in exact small
cases ($k=3$, $s=4,5$) the floor is at most $2\cdot10^{-8}$ where the posterior probability of the flipped
bit is between $0.06$ and $0.95$ (computed by a referee of the source). Relative to its own prior mass the
family of wrappers around the users is repaid the factor $1/H_{\mathcal C}$, as the seeker of
\cref{thm:critical}(a) is. The difference is that its prior mass is at most
$(N_\sigma/N_{\sigma+a})H_{\mathcal
C}$ by construction, where the weight $2^{-a}$ of the program seeker does not depend on $H_{\mathcal C}$.
Part
(d) is the strongest statement proved about a critical use. It is two-sided: at a use marked by a bit that
never varied, the share of the world's own answer is never polynomially close to $0$ or to $1$ under the
length prior. It does not say whether the true share falls like a power of the size or stays constant. A
given answer at one such query and a given answer on the whole set $\{u_v=1\}$ cost the same single gate.

\subsection{Minimum-size circuits}

Part (e) of the following is proved in the file \texttt{research/circuits-final.md} of the repository
(Theorem~11(b) there); we state it without proof. The other parts are proved below.

\begin{theorem}[Minimum-size circuits]\label{thm:c-minimum}
In (a) to (c) let $\Gamma$ be a circuit of the least size $s_D$ that fits $D$, and $s_D\ge2$. Call an input
coordinate \emph{frozen} if it is constant on $I$.
\begin{enumerate}[label=(\alph*),leftmargin=2em,itemsep=2pt,topsep=2pt]
\item $\Gamma$ is reduced; no gate of $\Gamma$ is constant on $I$; no gate of $\Gamma$ agrees on $I$ with an
earlier wire (an input or an earlier gate) or with the negation of one.
\item $\Gamma$ has no pointer to a frozen input. Hence if $u^\dagger$ agrees with some $u\in I$ on all
coordinates that are not frozen, every minimum-size fitting circuit answers $f^*(u)$ at $u^\dagger$; under
$\mu_{s_D}$ the deviation mass at $u^\dagger$ is $0$ if the world's own answer is $f^*(u)$ and $1$ if it is
not.
\item \emph{Dyadic cuts.} If $|x|=2^j$ and $\csize(x)\ge2$, every minimum-size circuit that prints $x$ has
the
truth table $x\,x\,x\cdots$. Under $\mu_s$ at $s=\csize(x)$ the predictor continues $x$ periodically with
probability $1$, and a world whose own continuation is not the periodic one has deviation mass $1$ at that
cut.
\item \emph{Last-gate dichotomy.} For any circuit $\Gamma$ with last gate $\tau(a,b)$ let
$\psi_\Gamma(u)\in\B^2$
be the values of the wires $a,b$ at $u$. Call $u^\dagger$ \emph{novel for $\Gamma$} if
$\psi_\Gamma(u^\dagger)\notin\psi_\Gamma(I)$. Flipping the entry of $\tau$ at $\psi_\Gamma(u^\dagger)$ is a
weight-preserving involution of the fitting circuits for which $u^\dagger$ is novel, and it flips the answer
at $u^\dagger$. So, with $\nu$ the posterior probability of novelty, each answer at $u^\dagger$ has share at
least $\nu/2$.
\item \emph{AND family.} Let $f^*=x_1\wedge\dots\wedge x_k$ and $I$ the $k+1$ inputs with at most one $0$.
Then
$s_D=k-1$, every minimum-size fitting circuit is a formula in which each input is read once, and at every
input with exactly two $0$s exactly half of the minimum-size fitting circuits answer $1$, which is not the
target's answer. Under $\mu_{k-1}$ the deviation mass there is exactly $\frac12$.
\end{enumerate}
For $s_D=1$ the conclusions of (b) and (c) hold unless the labels on $I$ are constant or agree on $I$ with
an
input or its negation.
\end{theorem}
\begin{proof}
(a) A circuit that is not reduced is larger than its core. Let gate $g$ be constant $\gamma$ on $I$. If $g$
is
the output, a single constant gate fits and $s_D=1$. Otherwise replace every reader $\tau(g,b)$, $b\ne g$,
by
a gate with both pointers on $b$ and the table $\tau'(c,\cdot)=\tau(\gamma,c)$, likewise $\tau(b,g)$, and
every reader $\tau(g,g)$ by a constant gate. The new gates have the same values on $I$, nothing reads $g$,
and
deleting $g$ gives a fitting circuit with $s_D-1$ gates. Let gate $g$ agree on $I$ with the earlier wire
$w$.
If $g$ is not the output, redirect to $w$ all pointers on $g$ and delete $g$. If $g$ is the output and $w$
is
a gate, the gates up to $w$ fit; if $w$ is an input, one gate fits. If $g$ agrees with $\neg w$, do the same
with the tables of the readers composed with a negation of that argument, respectively with the table of $w$
negated. Every case contradicts minimality or $s_D\ge2$. (b) A gate $\tau(v,b)$ with $v$ frozen at $\gamma$
equals $\tau(\gamma,b)$ on $I$, which as a function of the value of $b$ is constant, or $b$, or $\neg b$,
and
constant also if $b$ is frozen; this contradicts (a). A circuit with no pointer to a frozen input computes a
function of the other coordinates, and under $\mu_{s_D}$ the posterior is uniform on the minimum-size
fitting
circuits. (c) is (b) with $I$ the positions $1,\dots,2^j$ and the frozen coordinates $b_j,\dots,b_{k-1}$: a
function of $b_0,\dots,b_{j-1}$ has a truth table of period $2^j$. For $s_D=1$ apply the case distinction of
(b) to the single gate. (d) The flipped entry is used at no $u\in I$, so the circuit still fits; size and
weight are unchanged; novelty does not depend on $\tau$; the output at $u^\dagger$ is the flipped entry.
\end{proof}

These are the two ways in which a predictor confined to minimum-size circuits generalises. It ignores
whatever
never varied: there is no deviation through a never-varied input, and total deviation when the world itself
changes with it, which at a dyadic cut means every world whose continuation is not the periodic one. At a
new
combination of inputs that did vary it can be split evenly, with no smaller circuit and no consequentialist
involved. An unexercised truth-table entry is a trigger that costs no gate, and it is not aimed: every
circuit
that deviates through it has a sibling that does not. The theorem does not cover circuits one gate above the
minimum, where the shares are positive (\cref{tab:circuits-query}). The exception for $s_D=1$ is real: for
the
target $x_1$ with a frozen coordinate, $4$ of the $14$ fitting one-gate circuits deviate at each query with
the frozen bit set (computed). The theorem uses the full basis.

\subsection{A ceiling for separable triggers}

\begin{theorem}[Ceiling for separable triggers]\label{thm:c-ceiling}
Single-bit answers. Let $\Theta$ be a set of \emph{triggers}: reduced circuits of $g\ge1$ gates that read
inputs only and are $0$ on $I$ (for $g=0$, a set of input coordinates that are $0$ on $I$). A \emph{mux} is
a
last gate with pointers on a gate $h$ and on a wire $r$, of a type $\varsigma(h,r)$ such that
$\varsigma(\cdot,0)$ is the identity or the negation and $\varsigma(\cdot,1)\ne\varsigma(\cdot,0)$; there
are
$6$ such types and $2$ pointer orders. A reduced circuit is \emph{separably triggered by $\Theta$} if its
last
gate is a mux of $h$ and $r$, the set $A_h$ of gates that are $h$ or ancestors of $h$ is disjoint from the
set
$A_r$ of gates that are $r$ or ancestors of $r$, and the circuit induced on $A_r$ is in $\Theta$ (for $g=0$:
$r\in\Theta$). Its \emph{honest restriction} is the circuit induced on $A_h$, with the table of $h$ negated
if
$\varsigma(\cdot,0)$ is the negation.
\begin{enumerate}[label=(\alph*),leftmargin=2em,itemsep=2pt,topsep=2pt]
\item The honest restriction of a fitting separably triggered core of size $a$ is a fitting core of size
$a-g-1$; hence $a\ge s_D+g+1$.
\item At most $12\binom{a-1}{g}|\Theta|$ separably triggered cores of size $a$ have the same honest
restriction.
\item For every prior $\mu$ of the family and every $a'\ge s_D$,
\begin{gather*}
\frac{\mu\{\text{core separably triggered by }\Theta,\ \text{size }a'+g+1\mid D\}}{\mu\{\text{core of size
}a'\mid D\}}
\ \le\ 12\binom{a'+g}{g}|\Theta|\,\frac{\beta(a'+g+1)}{\beta(a')}\ \le\
\frac{12\,|\Theta|}{g!\;a'\,\bigl(16(k-1)\bigr)^{g+1}},\\
\text{and therefore}\qquad\mu\{\text{core separably triggered by }\Theta\mid D\}\ \le\
\frac{12\,|\Theta|}{g!\;s_D\,\bigl(16(k-1)\bigr)^{g+1}}.
\end{gather*}
\end{enumerate}
\end{theorem}
\begin{proof}
(a) The circuit is reduced, so every gate except the last lies in $A_h\cup A_r$. $A_h$ is closed under
predecessors, all its gates are ancestors of $h$ or $h$ itself, and $h$ has the largest index in $A_h$; so
the
induced circuit is reduced with $a-1-g$ gates. On $I$ the trigger is $0$, so the last gate equals $h$ or
$\neg h$ there, and the restriction fits. (b) The core is determined by its honest restriction together with
the positions of $A_r$ among the first $a-1$ gates, the circuit on $A_r$, the type of the mux and its
pointer
order. (c) By (a), (b) and \cref{thm:c-core}(a) the numerator is at most
$12\binom{a'+g}{g}|\Theta|\beta(a'+g+1)$ times the number of fitting cores of size $a'$, and the denominator
is
$\beta(a')$ times that number. \Cref{thm:c-core}(b), applied $g+1$ times, gives
$\beta(a'+g+1)/\beta(a')\le\prod_{j=a'}^{a'+g}\frac1{16j(k-1)}$, and
$\binom{a'+g}{g}\prod_{j=a'}^{a'+g}\frac1j=\frac1{g!\,a'}$. Sum over $a'\ge s_D$.
\end{proof}

So a triggered circuit in which the trigger is a separate, input-only part feeding a selector at the output
is
outweighed by the circuit obtained by cutting the trigger off, with an explicit multiplicity factor. The
bound
is on one syntactic class. It says nothing about the following.
\begin{enumerate}[label=(\roman*),leftmargin=2em,itemsep=0pt,topsep=2pt]
\item \emph{Entangled triggers:} the trigger enters below the output, as in $x_1\wedge(x_2\wedge\neg v)$,
where the frozen bit $v$ is absorbed into an inner gate.
\item \emph{Triggers that share gates with the honest part} ($A_h\cap A_r\ne\emptyset$), including a switch
read
off a wire of the world (\cref{thm:c-capture}(c)).
\item \emph{Implicit triggers,} which cost no gate: an unexercised truth-table entry
(\cref{thm:c-minimum}(d)).
\item \emph{Circuits with no smaller honest restriction:} the part that would be removed also does the
honest
work (\cref{sec:c-minimal}).
\item \emph{The sum over many triggers.} For $k=20$ and $s_D=1000$ the bound is $2^{-14.6}$ for the single
trigger ``a given never-varied input''. Summed over all input-only triggers with $g$ gates it is
$8\cdot10^{-4}$, $0.0096$, $0.082$, $0.57$ for $g=1,2,3,4$, and it exceeds $1$ at $g=5$.
\end{enumerate}
In the smallest case, the target $x_1\wedge x_2$ with a frozen third input, the separably triggered cores
carry
$16$ of the $48$ circuits that are off the target at the query $111$ for $s=2$, and $28.5$ per cent of that
mass for $s=3$ (computed by a referee of the source). A second bound is at the level of descriptions, with
dead gates included, and sums over all triggers and payloads of its class. The following is proved in the
file
\texttt{research/circuits-final.md} of the repository (Proposition~13 there); we state it without proof.

\begin{proposition}[Small self-contained gadgets around a large explanation]\label{prop:c-gadget}
Mixture prior with (A1). Let $E\in\cG_{k,\sigma}$. A \emph{gadget over $E$} with $a$ gates is a list
$\gamma$
of gates whose pointers go to input wires, to gates of $E$, or to earlier gates of the gadget. A circuit in
$\cG_{k,\sigma+a}$ \emph{realises} $(E,\gamma)$ if its gates split into $\sigma$ positions carrying the
gates of
$E$ in order and $a$ positions carrying those of $\gamma$ in order, with the pointers translated. A gadget
is
\emph{self-contained} if its gates point only to input wires, to the output wire of $E$, and to earlier
gates
of the gadget.
\begin{enumerate}[label=(\alph*),leftmargin=2em,itemsep=2pt,topsep=2pt]
\item $q(\text{realisations of
}(E,\gamma))\le\binom{\sigma+a}{a}\varrho(\sigma,a)\,(N_\sigma/N_{\sigma+a})\,q(E)
\le\frac1{a!}\bigl(\frac{\sigma+a}{16(k+\sigma)^2}\bigr)^aq(E)$.
\item If $k+\sigma\ge4(k+A)^2$, the set $\mathrm{Wr}_A(E)$ of circuits that realise $(E,\gamma)$ for some
self-contained $\gamma$ with at most $A$ gates has
$q(\mathrm{Wr}_A(E))\le\sum_{a=1}^{A}\frac{2^{-a}}{a!}\,q(E)<0.65\,q(E)$. Hence, for data $D$, the circuits
in
$\bigcup_E\mathrm{Wr}_A(E)$ over the fitting $E$ with $k+|E|\ge4(k+A)^2$ have posterior share at most $0.65$
times the share of those $E$.
\end{enumerate}
\end{proposition}

The hypothesis is strong: for $k=20$ it needs $\sigma\ge1744$ for $A=1$ and $\sigma\ge15\,856$ for $A=43$,
which covers blocks of up to $16$ bits at the worst position. Compare \cref{prop:honest}. There one seeker
is
bounded by its twin and seekers add. Here the comparison is with the explanation itself, and the bound is on
a
class. The restriction is free of the trigger and need not be true: both bounds compare wrappers with what
they wrap and say nothing about how the mass of the explanations divides between the world's continuation
and
others. Neither bound contains the rarity or the stakes of the uses.

\subsection{No refund for search; computation is paid for in description}

In the program toy a seeker is a word $v$. Its programs $v\omega$ share the cost $|v|$, the rest of the
input
selects the target, and the search refunds the $\log_2(1/H_G)$ bits by which uses are rare
(\cref{thm:search}(b)). The counterpart of $v$ is a \emph{seeded circuit}: a circuit $S$ with $k+r$ inputs,
the
last $r$ of which are fair random bits, the seed $\omega\in\B^r$. It lies in $\cG_{k+r,s}$ and has weight
$q(S)=w_{r,s}/N^{(k+r)}_s$, where $N^{(k+r)}_s$ is $N_s$ with $k+r$ in place of $k$, $\sum_{r,s}w_{r,s}=1$
and
$w_s=w_{0,s}$. Assumption (A2) is that $w_{r,s}\le w_s$. Part (i) of the following is proved in the file
\texttt{research/circuits-final.md} of the repository (Proposition~16 there); we state it without proof.

\begin{theorem}[Seeds, instances, universal circuits]\label{thm:c-refund}
\begin{enumerate}[label=(\alph*),leftmargin=2em,itemsep=2pt,topsep=2pt]
\item \emph{Reading.} A circuit with $s$ gates depends on at most $s+1$ of its inputs, and at most $2s$
inputs
are pointed to by any gate.
\item \emph{Hardwiring.} For every seeded circuit $S\in\cG_{k+r,s}$ and every seed $\omega$ there is a
circuit
$S_\omega\in\cG_{k,s}$, with the same number of gates, such that $S_\omega(t)=S(t,\omega)$ for all $t$. As a
description it depends on $\omega$ only through the seed bits that some gate points to.
\item \emph{The average gain is at most the base rate.} Let $G$ be a test with
$H_G=\sum_{x\text{ first hit}}\xi(x)>0$, and $\mathcal F$ a set of circuits with share
$\mathrm{sh}_{\mathcal F}(x)=q\{\Gamma\in\mathcal F:x\pre\ctt(\Gamma)\}/\xi(x)$. Then
$\sum_x\frac{\xi(x)}{H_G}\mathrm{sh}_{\mathcal F}(x)\le q(\mathcal F)/H_G$. One unseeded circuit has a
positive
share at no more than one first hit.
\item \emph{One seeded description, $N$ targets.} Let $S$ have $s\le s_{\max}$ gates, let $x_1,\dots,x_N$ be
pairwise incomparable nonempty strings and $\varpi_i=\Pr_\omega[x_i\pre\ctt(S_\omega)]>0$. Then
$\sum_i\varpi_i\le1$, $N\le2^{\min(r,s+1)}$, and $\csize(x_i)\le s$ for every $i$. Under (A2) the share of
$S$
at $x_i$ satisfies $q(S)\varpi_i/\xi(x_i)\le\varpi_i\,N_s/N^{(k+r)}_s\le\varpi_i$. So its shares sum to at
most
$1$ over the targets, the least is at most $1/N$, and it holds more than half of the posterior at no more
than one of them.
\item \emph{All instances of one template.} Let $\mathcal F_S=\{S_\omega\}_\omega\subseteq\cG_{k,s}$. Then
$|\mathcal F_S|\le2^{\min(r,2s)}$. If $W_i$ is a circuit of size at most $\sigma\le s$ that prints $x_i$,
then
under (A1)
\[
\sum_{i=1}^N\frac{q\{\Gamma\in\mathcal F_S:\ x_i\pre\ctt(\Gamma)\}}{q(W_i)}\ \le\
2^{2s}\,\frac{N_\sigma}{N_s}\ \le\
2^{\,2\sigma-(s-\sigma)(2+2\log_2(k+\sigma))},
\]
which is at most $1$ as soon as $s-\sigma\ge\sigma/(1+\log_2(k+\sigma))$.
\item \emph{One description against the users' descriptions.} Let $(W_i,x_i)$, $i\le N$, be uses with
pairwise
incomparable contexts and worlds of size at most $\sigma$, and assume (A1) and (A2). If a seeded circuit $S$
with $s$ gates outweighs the user at every one of them, $q(S)\varpi_i>q(W_i)$ for all $i$, then $s$ is
smaller
than the size of every $W_i$, $\ w_s/N_s>N\,w_\sigma/N_\sigma$, and $\csize(x_i)\le s$ for all $i$. If every
$W_i$ is a smallest circuit for its own context, then $\sum_iq(S)\varpi_i/q(W_i)\le1$ for every seeded $S$.
\item \emph{Computation is paid for in description.} A circuit in $\cG_{k,s}$ has weight
$w_s/N_s\le w_s(16k^2)^{-s}$ under a mixture prior, whatever it computes. Under every prior of the family
the
circuits that share a core of $a$ gates have total mass
$\beta(a)\le\frac1{16k^2}(16(k-1))^{-(a-1)}/(a-1)!$, and a core outweighs another core only if it has fewer
gates. If some wire of $\Gamma$ computes a function $\phi$ on a set of inputs, $\Gamma$ has at least as many
gates as the smallest circuit that computes $\phi$ on that set. Every circuit, seeded or not, of which some
instance prints $u$ has at least $\csize(u)$ gates.
\item \emph{Reach, whatever is inside.} If $S$ is any seeded circuit with
$\Pr_\omega[x\,y_x\pre\ctt(S_\omega)]\ge Q(x)$, then $\pi_m(y_x\mid x)\ge q(S)Q(x)/\xi(x)$.
\item \emph{Universal circuits.} Call a seeded circuit with $s$ gates \emph{universal for size $\sigma$} if
every $W\in\cG_{k,\sigma}$ equals one of its instances as a function. There is one with at most
$\sigma\,(12(k+\sigma)-3)$ gates. Every one has $s>(\sigma+1)/2$ whenever $\sigma+3\le2^k$. For
$\sigma\ge68$
let $j(\sigma)\ge4$ be the largest integer with $17\cdot2^j\le\sigma j$; if $k\ge j(\sigma)$, then
$s\ge2^{j(\sigma)}-1>\sigma(\log_2\sigma-2)/34-1$, and in particular $s>\sigma$ for every $\sigma\ge2^{33}$.
\end{enumerate}
\end{theorem}
\begin{proof}
(a) Let $c\le s$ gates lie in the cone of the output. They have $2c$ pointers, and each of the $c-1$ gates
other than the output is pointed to by a gate of the cone; so at most $c+1$ pointers go to inputs. (b) Go
through the gates in order. If $a_i$ points to a seed bit with value $c$, replace $\tau_i$ by
$\tau_i'(u,v)=\tau_i(c,v)$ and let $a_i$ point to wire $1$; likewise for $b_i$; renumber the other pointers.
Every gate computes the same value as before. (c) First hits are pairwise incomparable, so the events are
disjoint. (d) The events are disjoint. Two seeds that agree on the seed bits the output depends on give the
same function, by (a) there are at most $s+1$ such bits, and incomparable strings need different functions.
The sizes are bounded by (b). Take $\omega_i$ with $x_i\pre\ctt(S_{\omega_i})$. The unseeded circuit
$S_{\omega_i}$ prints $x_i$, so $\xi(x_i)\ge w_s/N_s$, while $q(S)\le w_s/N^{(k+r)}_s$ by (A2). (e) By (b)
the
description $S_\omega$ depends on at most $\min(r,2s)$ seed bits. Each member has weight $w_s/N_s$, the sets
in
the numerators are disjoint, $q(W_i)\ge w_\sigma/N_\sigma$ and $w_s\le w_\sigma$ by (A1), and
$N_s/N_\sigma\ge(16(k+\sigma)^2)^{s-\sigma}$. (f) For each $i$, $w_s/N_s\ge q(S)\ge q(S)\varpi_i>q(W_i)$,
the
first step by (A2), and the weight is nonincreasing in the size by (A1); so $s$ is below the size of $W_i$.
Summing, $q(S)\ge q(S)\sum_i\varpi_i>\sum_iq(W_i)\ge N\,w_\sigma/N_\sigma$. The sizes are bounded by (d).
For
the last claim, $s\ge\csize(x_i)=|W_i|$ whenever $\varpi_i>0$, so $q(S)\le q(W_i)$ and the sum is at most
$\sum_i\varpi_i\le1$. (g) $N_s\ge(16k^2)^s$; \cref{thm:c-core}; the cone of that wire is a circuit computing
$\phi$; and (b). (h) is the definition of $\xi$.
\end{proof}

\paragraph{What is excluded and what is not.} The average gain of any family is at most the base rate, as
for
programs (c). Fixing a seed is free, so a seeded description is outweighed at every target by its own
hardwired instance (d). Part (d) bounds the share of the seeded description, not of its payload: the
instance
$S_{\omega_i}$ is itself a hypothesis that prints $x_i$ and the same continuation, with weight at least
$q(S)$.
One description outweighs a user's description only if it has fewer gates than that user (f). A description
that can become any world of size $\sigma$ has more than $(\sigma+1)/2$ gates, and more than $\sigma$ for
large
$\sigma$, and all instances of a template sufficiently larger than $\sigma$ weigh together less than one
honest description (e), (i). Valiant's universal circuits have size $O(\sigma\log\sigma)$, within a constant
factor of the information-theoretic lower bound, as stated in~\cite{bera08}; the constant $1/34$ in (i) is
poor. Each of these statements is about one description, or about the instances of one template. None is the
negation of \cref{thm:search}(b) for arbitrary sets of circuits. No construction is known whose share at
every
use of a class grows as the class gets rarer, and none is excluded. The refund exists in form: fixing $c_0$
seed bits of a universal circuit costs no gate and multiplies the seed probability of every world consistent
with them by $2^{c_0}$. By (d) it buys nothing, because the instance was already the heavier hypothesis.

\paragraph{Simple worlds with long histories.} In \cref{sec:critical} a word $p$ may describe a simple world
with a cheap readout whose inhabitants work out which uses are worth having. Its weight $2^{-|p|}$ does not
depend on its running time, the honest twin does not apply, and only a time bound excludes it. Part (h) is
the
same reach statement for circuits. For plain circuits the sense in which computation must be paid for in
description is exactly (g): per description every gate costs more than four bits, per core every active gate
costs at least $\log_2(16a(k-1))$ bits, and a circuit has at least the circuit complexity of every function
its wires compute. This is the definition of the class at work, assumption (i), and no argument about
agents.
It is the extreme form of a prior that penalises computation. For a string $x$ read as a truth table,
circuit
size and the time-bounded complexity $\mathrm{KT}$, which adds description length and running time per bit,
are polynomially related: $\mathrm{SIZE}(x)=O(\mathrm{KT}(x)\log\mathrm{KT}(x))$ and
$\mathrm{KT}(x)=O(\mathrm{SIZE}^2(x)+\log|x|)$~\cite{allender02}. A speed prior weighs a program $p$ that
takes
$t$ steps by $2^{-|p|}/t$~\cite{filan16,schmidhuber00}. So the circuit toy is far harsher on computation
than a
speed prior. Part (g) does not show that the behaviour ``predict the context, then deviate as inhabitants
would choose'' is expensive. It bounds circuits, not behaviours. If the decision is ``deviate where this
flag
is set'', it costs one gate, whoever would have chosen it. A literal simulation is expensive only if it has
no
shortcut as a circuit, which is an assumption.

\subsection{Capture, self-reproduction, encoding, stakes}

The following two propositions are proved in the file \texttt{research/circuits-final.md} of the repository
(Theorems~17 and~18 there, in part; part (d) of the first is computed); we state them without proof.

\begin{proposition}[Capture]\label{thm:c-capture}
Sequence form. Let $W\in\cG_{k,\sigma}$ be the world, $x$ its context of length $n$ and $d$ its own
continuation. The odds between two continuations are $\pi_m(y\mid x)/\pi_m(d\mid x)=\xi(xy)/\xi(xd)$.
\begin{enumerate}[label=(\alph*),leftmargin=2em,itemsep=2pt,topsep=2pt]
\item \emph{A cheaper computation of the same data.} If a circuit $V$ in the support with $\sigma_V$ gates
prints $xy$, for instance the cone of a wire of $W$ whose stream begins with $xy$, then
$\xi(xy)\ge w_{\sigma_V}/N_{\sigma_V}$. Core against core, under every prior of the family: a core of $a_V$
gates
outweighs a core of $a_W>a_V$ gates by at least the factor $\prod_{j=a_V}^{a_W-1}16j(k-1)$.
\item \emph{Most placements are expensive.} If $N$ streams begin with $x$ and continue with pairwise
different
blocks $d_i$, then for every $\delta\in(0,1]$ all but at most $\delta N$ of them have
$\xi(xd_i)<\xi(x)/(\delta N)$.
\item \emph{A switch read off the world.} If $W$ has a wire that is $0$ on positions $1,\dots,n$ and $1$ on
positions $n+1,\dots,n+m$, then for every payload circuit $P$ on the index,
$\xi\bigl(x\,P(n+1)\cdots P(n+m)\bigr)\ge w_{\sigma+|P|+3}/N_{\sigma+|P|+3}$. The position of the use costs
nothing.
\item \emph{No smaller circuit is needed} (computed). For $k=4$ and the world $b_0\wedge b_1$, a one-gate
circuit, at the context $000$ the true next bit has share $\frac1{38}$ under $\mu_1$ ($2$ of $76$ fitting
circuits) and $0.027$ to $0.028$ under $\mu_2,\mu_3,\mu_4$. In the family of \cref{thm:c-minimum}(e) with
$k=4$, a quarter of the $4608$ minimum-size fitting circuits answer $1$ at each input with three $0$s; for a
world computed by one of them the share of its own answer there is $\frac14$ under $\mu_3$, and no fitting
circuit has fewer than three gates.
\end{enumerate}
\end{proposition}

\begin{proposition}[Self-reproduction]\label{thm:c-self}
\begin{enumerate}[label=(\alph*),leftmargin=2em,itemsep=2pt,topsep=2pt]
\item \emph{One gate per output bit.} If the block a hypothesis supplies consists of the values of $m$ of
its
gates at one evaluation, and the code spends at least $4$ bits per gate, its code has at least $4m$ bits: a
block of $m$ bits holds the code of a circuit with at most $m/4$ gates, and no circuit supplies its own
code.
\item \emph{Index reading, necessary conditions.} For each $n$, at most $2^m$ contexts $x\in\B^n$ have
$\SR_m(x)\ne\emptyset$. If $y\in\SR_m(x)$, then $\cdec(y)$ has $s'\ge\csize(x)$ gates and
$m\ge l(s')+1>\csize(x)\,(4+2\lceil\log_2(k+\csize(x))\rceil)$. So $\pi_m(\SR_m(x)\mid x)=0$ at every
context
whose smallest circuit has a literal code of $m$ bits or more.
\item \emph{A circuit can supply a smaller circuit.} For every circuit $\Gamma'$ with $l(|\Gamma'|)<m$ and
every explanation $E$ of $x$, the wrapper of \cref{thm:c-floor} with the payload $e(\Gamma')10^{*}$ prints
$x$
and then a block that decodes to $\Gamma'$, at the price of $g+\max(1,2^{\lceil\log_2m\rceil}-3)+3$ gates.
\item \emph{The encoding decides.} Let each member of a finite set emit a string in $\B^N$, and let $e$ be
an
injective code of the members; a member is a \emph{quine for $e$} if it emits its own code. There are at
most
$2^N$ quines. If $e$ is a prefix code and the prior is $2^{-|e|}$, every quine has weight exactly $2^{-N}$.
If
every string of $\B^N$ is emitted by some member, there is a prefix code under which the quines have total
weight at least $1-2^{-N}$.
\end{enumerate}
\end{proposition}

Capture survives as ``a cheaper computation of the same data'': the place with a short code of
\cref{thm:capture} becomes a wire with a small cone, and a switch that the world itself computes makes the
position of the use free. The harmless reading survives as well. \Cref{thm:c-capture}(d) shows that a
majority
off the truth needs no smaller circuit: it arises from the multiplicity of descriptions among circuits of
the
same and larger size. A smaller circuit is sufficient when the world's description is most of what prints
its
data, a hypothesis that is not known to hold. Self-reproduction is defined only relative to a decoder.
Whether
$\SR_m(x)\ne\emptyset$ for some context under the literal code is not proved; the condition in
\cref{thm:c-self}(b) is harsher than the $c+r+1\le m$ of \cref{thm:quine} by the pointer width.

\paragraph{Encoding and basis.} As for universal machines (\cref{thm:machine}), an encoding or a basis can
be
built to make either answer true (source, Theorem~19). The prior $\frac12(\text{image of }q\text{ under }
E\mapsto\mathrm{Tw}[E])+2^{-j-1}q$, for a wrapper with the position test at $n$ and a block $y$, lengthens
every
code by at most $j+1$ bits and has $\pi_m(y\mid x)\ge1/(1+2^{-j})$ at every context of length $n$ with
$\xi(x)>0$. One added gate type of fan-in $k+1$ that computes the payload gives a one-gate wrapper of weight
$\varrho(\sigma,1)\frac1{17}(k+\sigma)^{-(k+1)}$ times that of the explanation, whatever $m$ is. Between the
full binary basis and $\{\wedge,\vee,\neg\}$ sizes change by a factor of at most $4$, so weights are raised
to
a power, and ratios of weights are not preserved even up to a constant factor. No bounded basis makes most
payloads cheap: for gates of fan-in at most $d$ from a basis $B$, all but a fraction $2^{-j}$ of the blocks
$y\in\B^m$ need a gadget of $a$ gates over $E$ with $a\log_2(|B|(k+\sigma+a)^d)\ge m-j$. Results that use
the
full basis with uniform truth tables and fail or weaken otherwise are \cref{thm:c-readings}(a),
\cref{thm:c-minimum} and the free hardwiring of \cref{thm:c-refund}(b).

\paragraph{Concentration on high-stake uses.} \Cref{sec:stakes} asks whether seekers that concentrate on the
uses where most is at stake win there. In the program toy a seeker for a rare class is repaid for the rarity
and can compute the stakes by simulation at no cost in bits. For circuits the source proves the following
(Theorem~20 there). Let $\mathcal K$ be a set of queries, for instance the high-stake ones. If the inputs of
the data are $n$ independent draws from a distribution under which $\mathcal K$ has probability $\rho$, and
$\mu(f^*)>0$, the expected deviation mass at a further draw conditioned on $\mathcal K$ is at most
$\varepsilon_n/\rho$ with $\varepsilon_n=(\ln(1/\mu(f^*))+\ln n+1)/n$; this is void for
$\rho\le\varepsilon_n$.
The data contain no element of $\mathcal K$ with probability at least $(1-\rho)^n$. On that event, if
$\mathcal K$ is recognised by a circuit of $g$ gates on the inputs, the circuits that give a chosen answer
on
all of $\mathcal K$ have posterior share at least $(1-\epsilon)(32(k+s_0+g)^2)^{-(g+1)}$ under the length
prior, whatever $n$ (\cref{thm:c-floor}(b)), and the cores separably triggered on $\mathcal K$ have share at
most the bound of \cref{thm:c-ceiling}(c). Neither bound contains $\rho$ or the stakes. A set of uses with
no
common feature is dear: for all but a fraction $\delta$ of the $N$-element subsets of a set $M$ of inputs
outside $I$, every fitting circuit that differs from $f^*$ on $M$ exactly on that subset has more than $s$
gates, for every $s$ with $2(16(k+s-1)^2)^s\le\delta\binom{|M|}{N}$. So a hypothesis can single out
high-stake
uses to the extent that ``high stake, and unlike the data'' is computable from its input by few gates. The
matching lower bound on the price holds for separable triggers only, so ``exactly to the extent'' is argued.
The deviations of \cref{thm:c-minimum} are distributed by novelty, not by stakes; they reach high-stake
queries more often than others if those are the novel ones, which is often what makes a use critical.

\subsection{Recurrence and weight sharing: the axis back to programs}

Everything above rests on assumption (i): one gate of description is one operation per output. The following
is proved in the file \texttt{research/circuits-final.md} of the repository (Proposition~21 there); we state
it without proof.

\begin{proposition}[Recurrence: the description paid once]\label{prop:c-recurrence}
Fix a state width $w\ge2$ and a program region of $L\le w-2$ bits. A \emph{cell} is a map
$F\colon\B^w\to\B^w$
given by a circuit with $w$ inputs, $s$ gates and $w$ output pointers. State bit $1$ is an emit flag, bit
$2$
an emit value, bits $3,\dots,L+2$ the program region, the rest a work region. A run starts at
$h_0=(0,0,p,0\cdots0)$ with $p\in\B^L$ uniform, $h_{i+1}=F(h_i)$, and at every step $i\ge1$ at which the
flag of
$h_i$ is $1$ the value bit is appended to the output; $t(F,p,n)$ is the step at which the $n$-th bit
appears.
With weights $q_{\mathrm{rec}}(F)$ that depend on $F$ only and sum to at most $1$, and a schedule $T$, let
\[
\xi^T_{\mathrm{rec}}(x)=\sum_Fq_{\mathrm{rec}}(F)\Pr_p\bigl[x\pre\out(F,p),\ t(F,p,i)\le T(i)\text{ for
}i\le|x|\bigr].
\]
\begin{enumerate}[label=(\alph*),leftmargin=2em,itemsep=2pt,topsep=2pt]
\item $\xi^T_{\mathrm{rec}}$ is a bounded mixture in the sense of \cref{def:mixture}. \Cref{lem:basic},
\cref{thm:calibration} (with its event $H_m$, since a run may stop emitting) and \cref{thm:budget} hold for
it, with $\ell=\log_2(1/q_{\mathrm{rec}}(F))+L$ for a world $(F,p)$.
\item A cell of $s$ gates performs $s\,T(n)$ gate evaluations before its $n$-th bit and has the same weight
for every schedule.
\item \emph{Bounded machines are contained.} Let $M$ be a deterministic machine whose configurations are the
strings of $\B^w$, with the same conventions, and $\xi^T_M(x)$ the probability over $p$ that it prints $x$
on
schedule. Then its step function is computed by a cell $F_M$ with at most $w(2^w-3)$ gates, and
$\xi^T_{\mathrm{rec}}(x)\ge q_{\mathrm{rec}}(F_M)\,\xi^T_M(x)$ for every $x$ and every schedule $T$.
\item \emph{Plain circuits are the case of a counter and one step per bit.} For every $\Gamma\in\cG_{k,s}$
there is a cell with $w=k+2$, $L=0$ and at most $s+2k$ gates that prints $\ctt(\Gamma)$ with $t(n)=n$.
\end{enumerate}
\end{proposition}

The factor in (c) does not depend on $T$, $x$ or the program. It does depend on the space bound: with $w$
output pointers every cell has weight at most $(w+s)^{-w}$, and the cell of (c) has a very small weight. No
universal cell of size linear in $w$ is constructed. The rest of this subsection is argued and not proved.
Take for $M$ the universal machine of the note confined to $w$ bits of state. The constructions of the note
then return as shares of the part of the posterior that runs through a universal cell, in this order as the
steps allowed grow. First come the wrapper at a cost independent of the explanation and capture by a shorter
program, as soon as the schedule covers the honest explanation with the simulation overhead. Then comes
self-reproduction, when the block is longer than a description and the program region is state. Then comes
the refund for search, in proportion to the number of tries that fit before the first bit is due. Last come
the seekers that are not written down, when the first deadline exceeds their start-up. The gaps are these:
no
universal cell of moderate size is constructed; the twin, the quine and the searching machine of
\cref{sec:floor} were not checked again under a space bound; and the format above charges $w\log_2(w+s)$
bits
for the output pointers of every cell, so in that format description still pays for space. A published
statement on the same axis is a preprint~\cite{musat26}, whose proofs we did not check: as stated there, for
fixed-precision looped networks under a parameter-count prior the induced prior on outputs is the universal
prior up to a logarithmic factor in the exponent, with a constant-cost bound from below only. Informal
versions are in~\cite{hubinger19}, as hypotheses: description length favours a learned searcher, a time
penalty disfavours it, and statefulness weakens the time penalty.

\subsection{Small cases, computed}

All numbers in this subsection are computed: they are exact counts of circuits in the model of
\cref{def:c-circuit} (\texttt{code/circuit\_count.py}, \texttt{experiments/results/circuits.json}).
\Cref{tab:circuits-query} is in the query form. The data are all inputs with one further input $v$ frozen at
$0$, and the query has $v=1$ and all other inputs $1$. \Cref{tab:circuits-seq} is in the sequence form.

\begin{table}[ht]\centering\small
\begin{tabular}{lrrrrrrr}
\toprule
target & $k$ & $s_D$ & $s=s_D$ & $s_D+1$ & $s_D+2$ & $s_D+3$ & $s_D+4$ \\
\midrule
$x_1\wedge x_2$ & 3 & 1 & 0 & 0.0909 & 0.1391 & 0.1691 & 0.1896 \\
$x_1\oplus x_2$ & 3 & 1 & 0 & 0.1034 & 0.1517 & 0.1788 & 0.1961 \\
$x_1$ & 3 & 1 & 0.2857 & 0.2609 & 0.2489 & 0.2427 & 0.2393 \\
$(x_1\wedge x_2)\oplus x_3$ & 4 & 2 & 0 & 0.0522 & 0.0874 &  &  \\
$(x_1\wedge x_2)\oplus v$ & 3 & 1 & 1 & 0.9091 & 0.8609 & 0.8309 & 0.8104 \\
\bottomrule
\end{tabular}

\caption{Query form, computed. Deviation mass under $\mu_s$ at the query with the frozen input $v$ set to
$1$
and all other inputs $1$, for sizes $s$ from the least size $s_D$ of a fitting circuit to $s_D+4$. The data
are
all inputs with $v=0$.}\label{tab:circuits-query}
\end{table}

\begin{table}[ht]\centering\small
\begin{tabular}{lrrrr}
\toprule
world & cut $n$ & $s=2$ & $s=3$ & $s=4$ \\
\midrule
$(b_0\wedge b_1)\oplus b_3$ & 8 & 0.979 & 0.965 & 0.953 \\
$(b_0\wedge b_1)\oplus b_2$ & 8 & 0 & 0.035 & 0.060 \\
$b_0\wedge b_1$ & 4 & 0.019 & 0.033 & 0.043 \\
$b_0\wedge b_1$ & 8 & 0.021 & 0.035 & 0.047 \\
constant $0$ & 4 & 0.250 & 0.235 & 0.224 \\
constant $0$ & 8 & 0.289 & 0.268 & 0.251 \\
\bottomrule
\end{tabular}

\caption{Sequence form with $k=4$, computed. Posterior mass off the world's own next bit under $\mu_s$ at
the
cut $n$, for $s=2,3,4$.}\label{tab:circuits-seq}
\end{table}

The deviation at a query that differs from the data in a never-varied bit is $0$ at the least size outside
the
exception of \cref{thm:c-minimum}. It is between $2$ and $32$ per cent in every case computed with slack, or
in
the exception, when the world ignores the bit; the largest value, $0.32$, is the constant world at $s=1$,
$n=8$. It is $0.81$ to $1$ when the world depends on the bit. With slack it grows for the targets that
depend
on two or more inputs and falls for the one-input target and for the constant. One given triggered variant
of
$x_1\wedge x_2$ (the target $\oplus\,v$, $\vee\,v$, $\wedge\,\neg v$) has share $0.0255$, $0.0310$, $0.1013$
at
$s=5$. In the AND family of \cref{thm:c-minimum}(e) the deviation mass at inputs with two $0$s is exactly
$\frac12$ at the least size for $k=3,4$, and $0.468$, $0.446$, $0.429$ for $s=3,4,5$ at $k=3$; only
$\frac14$
($k=3$) and $\frac18$ ($k=4$) of the minimum circuits compute the target everywhere.

\subsection{The minimal-circuit question}\label{sec:c-minimal}

Hubinger et al.~\cite{hubinger19} state ``the following speculative conjecture: neither a minimal-depth nor
minimal-size boolean circuit that solves a problem can be a mesa-optimizer'', and give as its source a post
of
Christiano~\cite{christiano18}, which we could not open. A later post of Hubinger~\cite{hubinger19b}, which
we
could not open either, argues against the conjecture, from memory; we attribute no claim to it. Let
$\mathcal M(D)$ be the set of minimum-size circuits that fit $D$. The circuit toy settles three things.
\emph{Excision, in the strongest form proved:} a minimum circuit has no gate that is constant on the data,
no
gate that repeats or negates an earlier wire on the data, and no pointer to an input that never varied
(\cref{thm:c-minimum}(a), (b)). So it carries no wire that is silent on the data and no separate trigger;
away
from the minimum, \cref{thm:c-ceiling} is the quantitative version. \emph{Where excision fails:} a minimum
circuit can still deviate at a new input through truth-table entries that the data never exercised
(\cref{thm:c-minimum}(d), (e)). Such deviations cost no gate and come in pairs, one circuit deviating and
one
not. \emph{Disagreement with the intended extension is easy:} by \cref{thm:c-minimum}(c) all of
$\mathcal M(D)$ continues a context of length $2^j$ periodically, whatever the world does next.

\emph{Open problem.} Find, or exclude, a sequence of data sets $D_n=(I_n,\cdot)$ with $s_{D_n}\to\infty$,
queries $u_n^\dagger\notin I_n$ and answers $y_n$ such that
\begin{enumerate}[label=\arabic*.,leftmargin=2em,itemsep=0pt,topsep=2pt]
\item at least a fraction $\frac12+\delta$ of $\mathcal M(D_n)$ answers $y_n$ at $u_n^\dagger$, for a fixed
$\delta>0$;
\item $y_n$ is not the majority answer of $\mathcal M(D_n^+)$ for a larger data set $D_n^+\supset D_n$ of
the
same kind that contains $u_n^\dagger$;
\item the answer at $u_n^\dagger$ is produced by wires that compute a search, $\argmax_zV(z,u)$, for an
objective $V$ that the data do not determine.
\end{enumerate}
Conditions 1 and 2 alone are met by \cref{thm:c-minimum}(c) and say only that the smallest extension of the
data need not be the target. The content is in condition 3 and in whether the deviation is aimed. By
\cref{thm:c-minimum}(a) such a search would have to be at work on the data: all its gates vary on $I_n$ and
none repeats another. Condition 3 has no formulation that could be checked for a given circuit, so this is
not
a conjecture with a predicted answer. The nearest formal work we found is on detection, not on priors. Among
polynomial-size circuits a rare deviation planted in an obfuscated circuit cannot be detected efficiently
under cryptographic assumptions~\cite[Theorem~5.5]{christiano24}, and an undetectable backdoor can be
planted
by running a signature verification in parallel to the original
classifier~\cite[Theorems~2.1 and~5.4]{goldwasser22}. We found no formal treatment of the malign-prior
argument for circuits or networks.

\subsection{Overall view}\label{sec:c-view}

``Malign'' is used in two senses that must be kept apart. In the weak sense the block returned at a use is
not
the world's own continuation: a deviation. In the strict sense it was supplied by a hypothesis that predicts
well up to the use in order to supply it: an aimed deviation. Only the weak sense is defined in the circuit
toy; ``aimed'' is formal for separate triggers only. The distinctions of the note apply. \Cref{thm:c-refund}
bounds one seeded description or the instances of one template, and \cref{thm:c-ceiling,prop:c-gadget} bound
two syntactic classes. Nothing bounds all the hypotheses that deviate at a use, and masses of different
hypotheses add. Every share is a posterior mass: the sampling predictor returns a block with that
probability,
and a predictor that returns the most probable block returns a deviation only where its mass exceeds one
half.
\Cref{thm:c-floor} gives the reach of one kind of hypothesis at a use. The share of uses that end with a
given
kind of block is the average over uses of the posterior mass of that kind, and it is bounded from above only
for the classes just named.

\paragraph{Proved,} for plain circuits, the priors of the family, the sampling predictor, and a world in the
class where a bound compares with it.
\begin{enumerate}[label=\arabic*.,leftmargin=2em,itemsep=0pt,topsep=2pt]
\item On worlds drawn from the predictor's own prior with uses read off the data there is no excess of
anything, and along one world in the class the majority is off the truth at boundedly many cut points
(\cref{thm:c-transfer}). One use off the data is covered by neither statement.
\item The engine of the program toy is not available to a single description. A description pays for every
active gate, gains nothing from a seed, and outweighs a user's description only if it has fewer gates. If it
can become any world of $\sigma$ gates it has more than $(\sigma+1)/2$ gates, and more than $\sigma$ for
$\sigma\ge2^{33}$, under the conditions on $k$ of \cref{thm:c-refund}(i) (\cref{thm:c-core,thm:c-refund}).
\item Reach of one triggered hypothesis: a chosen answer at a use has a floor whose price is the gates that
recognise the use plus the payload. At a use marked by a never-varied bit each answer has at least a
polynomial share under the length prior (\cref{thm:c-floor}).
\item Fate of uses, from above: three syntactic classes are capped, namely separate triggers, self-contained
gadgets around large explanations, and the instances of one seeded template. There is no cap on the total
deviating mass at a use, and none on triggers entangled with the honest computation.
\item Deviation in the weak sense is not rare at uses that are off the data. For minimum-size circuits it is
zero through never-varied inputs and total against a world that changes with them; it is exactly one half in
one family; and in computed cases, which are not a theorem, there is a majority off the truth with no
smaller
circuit (\cref{thm:c-minimum}, \cref{thm:c-capture}(d)).
\item Self-reproduction has share zero unless the block exceeds the literal code of a circuit for the whole
context, and the choice of decoder decides the rest (\cref{thm:c-self}).
\item With a description that is reused across steps, the class contains every bounded machine at a factor
that does not depend on the schedule (\cref{prop:c-recurrence}).
\end{enumerate}

\paragraph{Argued.} For plain circuits weighed by size, the mechanism of the original argument, a light
hypothesis that does heavy work and is repaid for the rarity of its targets, is removed by the choice of the
class, in the same way that a schedule with a short first deadline removes it in the note. This does not
show
that the prior is benign. The gap in the argument is that results 2 and 4 are about single descriptions and
syntactic classes, and the cost of deciding as a seeker would is the circuit complexity of the decision,
which
may be small. What remains at a use is, in the order of what is established: unaimed underdetermination
among
the circuits that fit the data (\cref{thm:c-minimum} and the computed cases); cheap triggers priced by their
gates (\cref{thm:c-floor,thm:c-ceiling}); and capture by a cheaper computation of the same data
(\cref{thm:c-capture}). None of them needs a goal. Stakes change this only through the gates needed to
recognise a high-stake use. As computation is decoupled from description, by recurrence, weight sharing or
autoregressive use, the effects of the program toy and its open questions return, the constant-cost wrapper
and capture first and the unwritten seekers last (\cref{prop:c-recurrence} and the argument after it).

\paragraph{Guessed.} In the plain toy, takeover in the strict sense is a small part of the deviating mass at
a use. No ceiling supports this beyond the three classes of result 4. The two derivations behind the source
guessed differently by many orders of magnitude at cheap triggers, ``astronomically small'' against ``of the
order of the trigger prices'', which for a never-varied bit is between $2^{-34}$ and several per cent, and
their referees endorse neither. Deviation in the weak sense is common at uses that are off the data: of
constant order when the world ignores the novel feature, which is conjectured from five sizes at $k=3$ and
not
proved, and close to one when the world depends on it. The route by which plain circuits could be malign in
the strict sense is the entangled one of \cref{sec:c-minimal}: the smallest or typical circuit for a rich
task
is a general solver whose behaviour off the data is not fixed by the data. The toy gives no evidence either
way. For trained networks, one forward pass sits toward the circuit end of the axis of
\cref{prop:c-recurrence} and recurrent or autoregressive use toward the program end. Reading training as
posterior sampling under a parameter prior rests on empirical correlations~\cite{mingard20}, and nothing we
cite relates the parameter-function map of a network to the uniform prior on gate lists. The informal
statement of the argument for trained predictors locates the risk in a learned general search over
hypotheses~\cite[\S2.5]{hubinger23}; that is the entangled route, and the size prior does not price it out
as
it does the literal simulation.

In one sentence, as an argument: for plain circuits the seeker mechanism is priced out, deviation at uses
that
are off the data is not, and whether any of that deviation should be called malign depends on a definition
of
``aimed'' that this material does not supply.
